% Volume VI: Adversaries, Platforms, and Automated Judgment
% Part XI. Untrustworthy Actors
% Chapter 61: Strategic Transparency
% Spine: proof-spines/ch061-spine.md

\chapter{Strategic Transparency}
\label{ch:ch061}


\textbf{Strategic transparency} treats publicity as a field of action rather than a passive condition. Actors can exploit demands for openness by staging visible compliance, overwhelming auditors with excessive volume, disclosing information in unusable form, or moving the consequential activity into categories that monitoring systems do not inspect.

Visibility therefore cannot be equated with innocence. A disclosure regime may be perfectly public and still fail epistemically if what becomes visible is selected for theater rather than audit. The chapter's preface thus marks transparency as strategic terrain: one must ask not simply what has been shown, but what has been made visible, what has been buried under excess, and what has been displaced into the blind spots of the monitoring architecture.
